% 竖排书：\clearpage / \include 之后的下一页应排满整页，不能提前换页
% Issue #199: \include 多个文件时，从第二个文件开始显示异常
% https://github.com/open-guji/luatex-cn/issues/199
% 根因：自由模式按「本页已用宽度」决定换页，强制换页时没清零，
% 上一页的列宽被算进下一页，第二章首页只排了几列就换页。
\documentclass[twoside]{ltc-tw-vbook}
\setmainfont{TW-Kai}
\pageSetup{header-enabled=true, twoside=true}
\title{史記}

\begin{document}
\begin{正文}
\chapter{秦始皇本紀}

秦初并天下，令丞相、御史曰：「寡人以眇眇之身，興兵誅暴亂，賴宗廟之靈，六王咸伏其辜，天下大定。今名號不更，無以稱成功，傳後世。其議帝號。」丞相綰、御史大夫劫、廷尉斯等皆曰：「昔者五帝地方千里，其外侯服夷服，諸侯或朝或否，天子不能制。今陛下興
\clearpage
\chapter{項羽本紀}

秦初并天下，令丞相、御史曰：「寡人以眇眇之身，興兵誅暴亂，賴宗廟之靈，六王咸伏其辜，天下大定。今名號不更，無以稱成功，傳後世。其議帝號。」丞相綰、御史大夫劫、廷尉斯等皆曰：「昔者五帝地方千里，其外侯服夷服，諸侯或朝或否，天子不能制。今陛下興義兵，誅殘賊，平定天下，海內爲郡縣，法令由一統，自上古以來未嘗有，五帝所不及。臣等謹與博士議曰：『古有天皇，有地皇，有泰皇，泰皇最貴。』臣等昧死上尊號，王爲『泰皇』。命爲『制』，令爲『詔』，天子自稱曰『朕』。」王曰：「去『泰』，著『皇』，采上古『帝』位號，號曰『皇帝』。他如議。」制曰：「可。」追尊莊襄王爲太上皇，制曰：「朕聞太古有號毋諡，中古有號，死而以行爲諡。如此，則子議父，臣議君也，甚無謂，朕弗取焉。自今已來，除諡法。朕爲始皇帝。後世以計數，二世、三世至於萬世，傳之無窮。」

丞相綰等言：「諸侯初破，燕、齊、荊地遠，不爲置王，毋以填之。請立諸子，唯上幸許。」始皇下其議於羣臣，羣臣皆以爲便。廷尉李斯議曰：「周文、武所封子弟同姓甚衆，然後屬疏遠，相攻擊如仇讎，諸侯更相誅伐，周天子弗能禁止。今海內賴陛下神靈一統，皆爲郡縣，諸子功臣以公賦稅重賞賜之，甚足易制。天下無異意，則安寧之術也。置諸侯不便。」始皇曰：「天下共苦戰鬬不休，以有侯王。賴宗廟，天下初定，又復立國，是樹兵也，而求其寧息，豈不難哉？廷尉議是。」

分天下以爲三十六郡，郡置守、尉、監。更名民曰「黔首」。大酺。收天下兵，聚之咸陽，銷以爲鍾鐻，金人十二，重各千石，置廷宮中。一法度衡石丈尺，車同軌。書同文字。地東至海暨朝鮮，西至臨洮、羌中，南至北嚮戶，北據河爲塞，並陰山至遼東。徙天下豪富於咸陽十二萬戶。諸廟及章臺、上林皆在渭南。
\end{正文}
\end{document}
